\section{The Coincidence}

Two facts sit side by side in the recent history of the Catholic Church, and neither is seriously disputed. The first is that between 1962 and 1970 the Church held an ecumenical council and then, in its wake, replaced the most stable object in her daily life: the Roman Missal, the altar book whose 1570 edition Paul VI himself described as one of the ``numerous and admirable fruits'' that the Council of Trent ``has spread throughout the entire Church of Christ,'' and which had furnished the Latin Rite's norm of Eucharistic celebration for four centuries.\endnote{Paul VI, apostolic constitution \emph{Missale Romanum} (3 April 1969), opening paragraph, official English text on the Holy See portal, checked 2026-07-25; the promulgated Latin is authoritative and was checked the same day as reprinted in the registered reproduction of the 2002 third typical edition (see References). Here and throughout, ``the 1962 Missal'' means the \emph{Missale Romanum} in the typical edition promulgated by John XXIII in 1962, and ``the Missal of Paul VI'' means the postconciliar \emph{Missale Romanum} in its typical editions of 1970, 1975, and 2002 (with the 2008 emended reprint of the third edition); the edition identifications are documented in section 4.} The second is that in the same decade, across most of the Western world, nearly every measurable index of Catholic practice turned downward---some gently, some off a cliff---and half a century later most have not recovered. In the United States, the best long series available reports that 54.9 percent of self-identified Catholics attended Mass in a typical week in 1970; in 2025 the same series reports 17.6 percent.\endnote{Center for Applied Research in the Apostolate (CARA), \emph{Frequently Requested Church Statistics}, workbook state downloaded 2026-07-25 from the CARA site (see References), row ``Catholics who attend Mass every week (survey-based estimate).'' The series begins at 1970 and offers no 1965 figure---a limit kept in section 5. The percentages are survey estimates of self-identified Catholics, not counts.} Religious sisters in the same country numbered 178,740 in 1965 and 33,135 in 2025.\endnote{CARA, \emph{Frequently Requested Church Statistics}, row ``Religious sisters,'' checked 2026-07-25; \emph{Official Catholic Directory} reported counts, whose definitions and limits are stated in section 5.}

Put the two facts in one sentence and a syllogism forms almost by itself. Before the Council: full seminaries, full convents, full pews. Then the Council, and the new Mass. After: collapse. \emph{Post hoc, ergo propter hoc} is a famous fallacy, but it is famous precisely because the inference is so natural, and here the chronological neighbors are enormous. One party in the Church has drawn the inference plainly: the Council---or its liturgical constitution, or the Missal produced in its name---broke something, and the numbers are the sound of it breaking. Another party reverses the sign: the storm was coming anyway; the Council built the only ship that could sail in it; where the reform was betrayed or half-applied, the losses were worse. A third position, common in official documents, declines the causal question and speaks of a crisis of faith afflicting the whole modern West, within which council and Missal are responses rather than causes.

These are not three moods; they are three incompatible causal claims, and the dispute has never been merely academic: it has produced an episcopal society in formal rupture with Rome, two papal legislative frameworks in sixteen years pointing in opposite practical directions, and, in ordinary parishes, a low-grade argument any Catholic under seventy has overheard. Yet the striking thing, once one reads the primary documents instead of the polemics, is how much of the argument on every side proceeds from texts unread and numbers unexamined. The constitution on the liturgy is invoked constantly and quoted rarely. The Missal of Paul VI is treated as the Council's direct product by both its fiercest critics and some of its defenders, though the Council closed before the Missal existed. The statistics are recited in whichever truncation serves, stripped of definitions, geography, and inconvenient counter-trends.

So this essay does the pedestrian thing first and lets the argument earn its conclusions in order: what the Council was summoned and claimed authority to do; what the Constitution on the Sacred Liturgy actually mandates, permits, and forbids; by what acts the Missal came to be and how it stands to the mandate; what the crisis indicators genuinely show once their ceilings are stated; how the rival readings of the whole sequence---rupture and reform---were framed by the pope who posed the question most exactly; where the traditionalist critique is strong, where it fails, and what the Church has conceded to it; and how the present discipline, from the 1984 indult through \emph{Traditionis custodes} and its 2023 rescript, governs a disagreement it has not resolved. One warning is owed at the outset: on the causal question this essay will end with less than the partisans promise, because the evidence ends there too. What it will not do is blur the three objects of the title into one. A council, a missal, and a crisis are different kinds of thing; nothing true can be said about them while they are fused.
